% =============================================================
% 02_accounting.tex — 会計・財務と法務
% =============================================================
\chapter{会計・財務と法務}

\chaptermeta{%
  \item PLの5段階利益を読み取れる
  \item 損益分岐点（BEP）を計算できる
  \item 知的財産権の種類と保護期間を整理できる
  \item 請負・準委任・派遣の違いを即答できる
  \item 情報倫理・コンプライアンスの基本を押さえる
}{75分}{2}{第1章}

ストラテジ系の後半は「お金」と「法律」です。
PLの5段階利益、損益分岐点の計算、知的財産権、契約形態——
\textbf{数字の出る問題}と\textbf{法律の識別問題}を確実に得点源にしましょう。

% =============================================================
\section{財務三表の全体像}
% =============================================================

\begin{teigi}[財務三表とは]
企業の経営状態を把握するための3つの書類。
\begin{itemize}
  \item \termtag{損益計算書（PL）}：一定期間の「儲け」を示す（フロー）
  \item \termtag{貸借対照表（BS）}：ある時点の「財産」を示す（ストック）
  \item \termtag{キャッシュフロー計算書（CF）}：資金の流れを示す
\end{itemize}
Iパスでは特にPLの出題頻度が高い。BSは「資産＝負債＋純資産」を押さえればOK。
\end{teigi}

\begin{pitfall}[BSの左右を逆にしがち]
貸借対照表は\textbf{左側（借方）＝資産}、\textbf{右側（貸方）＝負債＋純資産}。
「左に持っているもの、右にその調達元」と覚える。
\end{pitfall}

% =============================================================
\section{PLの5段階利益}
% =============================================================

\begin{teigi}[PLの5段階利益]
PLは売上高から順に費用を引いて5つの利益を計算する。

\begin{enumerate}
  \item \termtag{売上総利益}＝ 売上高 $-$ 売上原価（粗利）
  \item \termtag{営業利益}＝ 売上総利益 $-$ 販管費（本業の儲け）
  \item \termtag{経常利益}＝ 営業利益 $+$ 営業外収益 $-$ 営業外費用
  \item \termtag{税引前当期純利益}＝ 経常利益 $+$ 特別利益 $-$ 特別損失
  \item \termtag{当期純利益}＝ 税引前当期純利益 $-$ 法人税等
\end{enumerate}
\end{teigi}

\begin{hikaku}[PLの5段階を一覧比較]
\comptable{利益 & 計算式 & 何がわかるか}{
売上総利益 & 売上高 $-$ 売上原価 & 商品力 \\
営業利益 & 売上総利益 $-$ 販管費 & 本業の収益力 \\
経常利益 & 営業利益 $\pm$ 営業外損益 & 通常の経営力 \\
税引前当期純利益 & 経常利益 $\pm$ 特別損益 & 臨時含む全体像 \\
当期純利益 & 税引前 $-$ 法人税等 & 最終的な儲け
}
\end{hikaku}

% --- 例題 ---
\begin{reidai}{PLの計算（当期純利益から販管費を逆算）}
ある企業のPLが以下のとおりであるとき、販売費及び一般管理費は何百万円か。

\begin{center}
\small
\begin{tabular}{lr}
\toprule
売上高 & 5,000百万円 \\
売上原価 & 3,000百万円 \\
営業外収益 & 100百万円 \\
営業外費用 & 200百万円 \\
特別利益 & 50百万円 \\
特別損失 & 50百万円 \\
法人税等 & 100百万円 \\
当期純利益 & 800百万円 \\
\bottomrule
\end{tabular}
\end{center}

\sentakushi{800}{900}{1,000}{1,100}
\end{reidai}

\begin{shishin}
当期純利益から逆算する。税引前当期純利益 $=$ 当期純利益 $+$ 法人税等 $= 800 + 100 = 900$。
経常利益 $=$ 税引前 $-$ 特別利益 $+$ 特別損失 $= 900 - 50 + 50 = 900$。
営業利益 $=$ 経常利益 $-$ 営業外収益 $+$ 営業外費用 $= 900 - 100 + 200 = 1{,}000$。
売上総利益 $=$ 売上高 $-$ 売上原価 $= 5{,}000 - 3{,}000 = 2{,}000$。
販管費 $=$ 売上総利益 $-$ 営業利益 $= 2{,}000 - 1{,}000 = 1{,}000$。
\end{shishin}

\begin{kaitou}
\textbf{正解：ウ（1,000百万円）}

5段階利益を下から逆算する典型パターン。
当期純利益→税引前→経常→営業の順に巻き戻し、最後に売上総利益との差で販管費を求める。
\end{kaitou}

\begin{minuki}[見抜き方]
\begin{itemize}
  \item PLの問題は「逆算」がほとんど——与えられた利益から上方向に戻して未知を求める
  \item 特別損益がゼロ（相殺）なら経常利益＝税引前当期純利益と気づけるとスピードアップ
  \item 営業外損益の符号を間違えやすい——「収益はプラス、費用はマイナス」を徹底
\end{itemize}
\end{minuki}

% =============================================================
\section{損益分岐点}
% =============================================================

\begin{teigi}[損益分岐点（BEP）]
利益がちょうどゼロになる売上高。
\[
  \text{BEP} = \frac{\text{固定費}}{1 - \text{変動費率}}
  = \frac{\text{固定費}}{\text{限界利益率}}
\]
\begin{itemize}
  \item \termtag{固定費}：売上に関係なく一定（家賃、人件費など）
  \item \termtag{変動費}：売上に比例して増減（材料費など）
  \item \termtag{変動費率}＝ 変動費 ÷ 売上高
  \item \termtag{限界利益率}＝ 1 $-$ 変動費率
\end{itemize}
\end{teigi}

\begin{pitfall}[固定費と変動費の分類ミス]
\begin{itemize}
  \item 人件費は\textbf{固定費}（パート・アルバイトの時給制でも基本は固定費扱い）
  \item 材料費・仕入原価は\textbf{変動費}
  \item 問題文の「売上に比例する」かどうかで判断
\end{itemize}
\end{pitfall}

% =============================================================
\section{知的財産権}
% =============================================================

\begin{hikaku}[知的財産権の比較]
\comptable{権利 & 保護対象 & 保護期間}{
特許権 & 発明（技術的アイデア） & 出願から20年 \\
実用新案権 & 考案（物品の形状等） & 出願から10年 \\
意匠権 & デザイン & 出願から25年 \\
商標権 & ブランド名・ロゴ & 登録から10年（更新可） \\
著作権 & 創作的な表現 & 著作者の死後70年
}
\end{hikaku}

\begin{teigi}[著作権の特徴]
\begin{itemize}
  \item \termtag{無方式主義}：登録不要で、創作した時点で自動的に発生
  \item プログラムも著作権で保護される（ただしアルゴリズムは保護されない）
  \item アイデアそのものは著作権では保護されない（表現が対象）
\end{itemize}
\end{teigi}

\begin{pitfall}[特許権と著作権の混同]
\begin{itemize}
  \item \textbf{特許権}：出願・審査が必要（方式主義）、技術的アイデアを保護
  \item \textbf{著作権}：登録不要（無方式主義）、表現を保護
  \item 「プログラムの著作権」と「ソフトウェア特許」は別物
\end{itemize}
\end{pitfall}

% =============================================================
\section{契約形態（請負・準委任・派遣）}
% =============================================================

\begin{hikaku}[3つの契約形態の比較]
\comptable{項目 & 請負 & 準委任 & 派遣}{
完成責任 & あり（成果物を納品） & なし（善管注意義務） & なし \\
指揮命令 & 受注者 & 受注者 & 派遣先 \\
瑕疵担保 & あり & なし & --- \\
報酬の対象 & 成果物 & 作業時間・工数 & 労働時間
}
\end{hikaku}

\begin{teigi}[偽装請負とは]
\termtag{偽装請負}：形式上は請負契約だが、実態は発注者が労働者に直接指揮命令している状態。
\textbf{違法}であり、労働者派遣法に違反する。
\begin{itemize}
  \item 請負なのに発注者が作業の進め方を細かく指示 → 偽装請負
  \item 指揮命令権が発注者にあるなら、正しくは「派遣契約」を結ぶべき
\end{itemize}
\end{teigi}

% =============================================================
\section{情報倫理とコンプライアンス}
% =============================================================

\begin{teigi}[主要な法律]
\begin{itemize}
  \item \termtag{個人情報保護法}：個人情報の取扱いルール（利用目的の特定・公表）
  \item \termtag{不正アクセス禁止法}：他人のID/パスワードの不正使用を禁止
  \item \termtag{特定電子メール法}：広告宣伝メールのオプトイン規制
  \item \termtag{不正競争防止法}：営業秘密の不正取得・ドメイン名の不正使用を禁止
\end{itemize}
\end{teigi}

% =============================================================
\section{過去問ドリル}
% =============================================================

\shoumatsuhead{過去問ドリル：会計・法務}

\begin{itpmon}[\sourcebadge{R7}\enspace\diffbadge{標}]{%
A社がシステム開発をB社に委託している。A社の社員がB社の作業者に対して、日々の作業の割振りや進捗管理を直接行っている場合、この状態を何というか。}
\sentakushi
  {SES契約}
  {偽装請負}
  {派遣契約}
  {準委任契約}
\end{itpmon}
\begin{kaisetsu}{イ}
請負や準委任では指揮命令権は受注者側にある。発注者が直接作業指示をしている実態があれば、形式に関わらず\textbf{偽装請負}にあたる。ウの派遣契約であれば合法だが、請負契約のまま指揮命令していることが問題。
\end{kaisetsu}

\begin{itpmon}[\sourcebadge{R7}\enspace\diffbadge{標}]{%
特定電子メール法で規制される行為として、最も適切なものはどれか。}
\sentakushi
  {企業が自社の社員全員に業務連絡のメールを送信した。}
  {受信者の同意を得ずに、広告宣伝の電子メールを送信した。}
  {知人に転職の挨拶メールを送信した。}
  {取引先に見積書をメールで送信した。}
\end{itpmon}
\begin{kaisetsu}{イ}
\textbf{特定電子メール法}は「広告宣伝メール」に対して\textbf{オプトイン規制}（事前同意なしの送信を禁止）を定めている。業務連絡や個人的な連絡は規制対象外。
\end{kaisetsu}

\begin{itpmon}[\sourcebadge{R6}\enspace\diffbadge{標}]{%
著作権に関する記述として、最も適切なものはどれか。}
\sentakushi
  {著作権を取得するには、特許庁に登録する必要がある。}
  {著作権は、著作物を創作した時点で自動的に発生する。}
  {アルゴリズムは著作権で保護される。}
  {著作権の保護期間は、著作者の死後100年である。}
\end{itpmon}
\begin{kaisetsu}{イ}
著作権は\textbf{無方式主義}——登録不要で創作時に自動発生する。アは特許権の手続き。ウについてアルゴリズム（手順）は著作権の対象外。エの保護期間は死後\textbf{70年}。
\end{kaisetsu}

% =============================================================
\section{タイムアタック}
% =============================================================

\begin{timeattack}{会計・法務クイックチェック}{5分}{4問中3問}

\begin{itpmon}[\diffbadge{基}]{PLにおいて「売上高 $-$ 売上原価」で求められる利益はどれか。}
\sentakushi{営業利益}{経常利益}{売上総利益}{当期純利益}
\end{itpmon}
\begin{kaisetsu}{ウ}
売上高 $-$ 売上原価 $=$ \textbf{売上総利益}（粗利）。営業利益はさらに販管費を引いたもの。
\end{kaisetsu}

\begin{itpmon}[\diffbadge{標}]{固定費が600万円、変動費率が40\%のとき、損益分岐点売上高は何万円か。}
\sentakushi{800}{1,000}{1,200}{1,500}
\end{itpmon}
\begin{kaisetsu}{イ}
BEP $= 600 \div (1 - 0.4) = 600 \div 0.6 = 1{,}000$万円。
\end{kaisetsu}

\begin{itpmon}[\diffbadge{基}]{出願から20年間保護される知的財産権はどれか。}
\sentakushi{意匠権}{実用新案権}{商標権}{特許権}
\end{itpmon}
\begin{kaisetsu}{エ}
\textbf{特許権}は出願から20年。意匠権は25年、実用新案権は10年、商標権は登録から10年（更新可）。
\end{kaisetsu}

\begin{itpmon}[\diffbadge{標}]{請負契約と準委任契約の最も大きな違いはどれか。}
\sentakushi
  {契約の相手方が法人か個人かの違い}
  {成果物の完成責任があるかないかの違い}
  {作業場所がどこかの違い}
  {報酬が月額か日額かの違い}
\end{itpmon}
\begin{kaisetsu}{イ}
請負は\textbf{成果物の完成責任}がある。準委任は善管注意義務はあるが完成責任はない。
\end{kaisetsu}

\end{timeattack}

% =============================================================
\section{よくあるミスと到達確認}
% =============================================================

\begin{yokumiss}[この章でよくあるミス]
\begin{enumerate}[leftmargin=1.6em, itemsep=5pt]
  \item \textbf{営業利益と経常利益を混同する}——営業利益は本業の儲け、経常利益は営業外損益を含む。
  \item \textbf{損益分岐点の公式で変動費率と限界利益率を逆にする}——分母は$1-$変動費率（＝限界利益率）。
  \item \textbf{特許権と著作権の取得方法を混同する}——特許は出願・審査が必要、著作権は自動発生。
  \item \textbf{請負と準委任の指揮命令権を間違える}——どちらも指揮命令は受注者。派遣だけが派遣先。
  \item \textbf{意匠権の保護期間を忘れる}——2020年改正で出願から25年に延長された。
\end{enumerate}
\end{yokumiss}

\begin{tatsaku}
  \item PLの5段階利益を順番に言える
  \item 損益分岐点の公式を使って計算できる
  \item 知的財産権5種の保護対象と期間を言える
  \item 著作権の無方式主義を説明できる
  \item 請負・準委任・派遣の3点比較表を再現できる
  \item 偽装請負の判断ポイントを説明できる
  \item 特定電子メール法のオプトイン規制を理解している
  \item BSの左右（資産＝負債＋純資産）を説明できる
\end{tatsaku}

\nextchapter{%
マネジメント系に入ります。システム開発の5工程、ウォーターフォールとアジャイルの比較、
プロジェクトマネジメントのQCD——「つくる側」の視点を学びます。
}
